%%%%%%%%%%%%%%%%%%%%%%%%%%%%%%%%%%%%%%%%%%%%%%%%%%%%%%%%%%%%%%%%%%%%%%%%%%%%
%% Trim Size: 9.75in x 6.5in
%% Text Area: 8in (include Runningheads) x 5in
%% ws-ijtaf.tex   :   29-07-2021
%% Tex file to use with ws-ijtaf.cls written in Latex2E.
%% The content, structure, format and layout of this style file is the
%% property of World Scientific Publishing Co. Pte. Ltd.
%% Copyright 2021 by World Scientific Publishing Co.
%% All rights are reserved.
%%%%%%%%%%%%%%%%%%%%%%%%%%%%%%%%%%%%%%%%%%%%%%%%%%%%%%%%%%%%%%%%%%%%%%%%%%%%
%%

\documentclass{ws-ijtaf}
\begin{document}

\markboth{Authors' Names}{Instructions for Typing Manuscripts (Paper's Title)}

%%%%%%%%%%%%%%%%%%%%% Publisher's Area please ignore %%%%%%%%%%%%%%%
\catchline{}{}{}{}{}
%%%%%%%%%%%%%%%%%%%%%%%%%%%%%%%%%%%%%%%%%%%%%%%%%%%%%%%%%%%%%%%%%%%%

\title{INSTRUCTIONS FOR TYPESETTING MANUSCRIPTS \\
USING \LaTeX\footnote{For the title, try not to
use more than 3 lines. Typeset the title in 10 pt bold and uppercase.}
}

\author{FIRST AUTHOR\footnote{
Typeset names in 8 pt roman, uppercase. Use the footnote to indicate the
present or permanent address of the author.}}

\address{University Department, University Name, Address\\
City, State ZIP/Zone,Country\,\footnote{State completely without
abbreviations, the affiliation and mailing address, including country.
Typeset in 8 pt italic.}\\
\email{author\_id@domain\_name\footnote{Typeset author e-mail
address in single line.}} }

\author{SECOND AUTHOR}

\address{Group, Laboratory, Address\\
City, State ZIP/Zone, Country\\
author\_id@domain\_name }

\maketitle

\begin{history}
\received{(Day Month Year)}
\revised{(Day Month Year)}
\end{history}

\begin{abstract}
The abstract should summarize the context, content and conclusions of
the paper in less than 200 words. It should not contain any reference
citations or displayed equations. Typeset the abstract in 8 pt roman
with baselineskip of 10 pt, making an indentation of 1.5 pica on the
left and right margins.
\end{abstract}

\keywords{Keyword1; keyword2; keyword3.}

\section{General Appearance}	
Contributions to the {\it International Journal of Theoretical and
Applied Finance} will mostly be processed by using the authors' source
files. These should be submitted with the manuscripts, and resubmitted
in the final form if a paper requires revision before being accepted
for publication.

\section{The Main Text}

Authors are encouraged to have their contribution checked for grammar.
American spelling should be used. Abbreviations are allowed but should
be spelt out in full when first used. Integers ten and below are to be
spelt out. Italicize foreign language phrases (e.g.~Latin, French).

The text should be in 10 pt roman, single spaced with
baselineskip of 13~pt. Text area (excluding copyright block and folio)
is 6.9 inches high and 5 inches wide for the first page.
Text area (excluding running title) is 7.7 inches high and\break
5 inches wide for subsequent pages.  Final pagination and
insertion of running titles will be done by the publisher.

\section{Major Headings}
Major headings should be typeset in boldface with the first letter of
important words capitalized.

\subsection{Sub-headings}
Sub-headings should be typeset in boldface italic and capitalize the
first letter of the first word only. Section number to be in boldface
roman.

\subsubsection{Sub-subheadings}
Typeset sub-subheadings in medium face italic and capitalize the first
letter of the first word only. Section numbers to be in roman.

\subsection{Numbering and spacing}
Sections, sub-sections and sub-subsections are numbered in Arabic.
Use double spacing before all section headings, and single spacing
after section headings. Flush left all paragraphs that follow after
section headings.

\subsection{Lists of items}
List may be presented with each item marked by bullets and numbers.

\subsection*{Bulleted items}

\begin{itemlist}
 \item item one,
 \item item two,
 \item item three.
\end{itemlist}

\subsection*{Numbered items}

\begin{arabiclist}
 \item item one,
 \item item two,
 \item item three.
\end{arabiclist}

The order of subdivisions of items in bullet and numbered lists
may be presented as follows:

\subsection*{Bulleted items}

\begin{itemize}
\item First item in the first level
\item Second item in the first level
\begin{itemize}
\item First item in the second level
\item Second item in the second level
\begin{itemize}
\item First item in the third level
\item Second item in the third level
\end{itemize}
\item Third item in the second level
\item Fourth item in the second level
\end{itemize}
\item Third item in the first level
\item Fourth item in the first level
\end{itemize}

\subsection*{Numbered items}

\begin{arabiclist}
\item First item in the first level
\item Second item in the first level
\begin{alphlist}[(a)]
\item First item in the second level
\item Second item in the second level
\begin{romanlist}[iii.]
\item First item in the third level
\item Second item in the third level
\item Third item in the third level
\end{romanlist}
\item Third item in the second level
\item Fourth item in the second level
\end{alphlist}
\item Third item in the first level
\item Fourth item in the first level
\end{arabiclist}

\section{Equations}

Displayed equations should be numbered consecutively, with the number
set flush right and enclosed in parentheses. The equation numbers
should be consecutive within each section.  Numbering by sections
refers only to the large main sections, i.e. to those with Level 1
headings. Thus the equations in Sec.~4 would be labeled (4.1),
(4.2), etc.
\begin{equation}
\mu(n, t) = \frac{\sum^\infty_{i=1} 1(d_i < t, N(d_i)
= n)}{\int^t_{\sigma=0} 1(N(\sigma) = n)d\sigma}\,.
\label{eq:jaa}
\end{equation}

Equations should be referred to in abbreviated form,
e.g.~``Eq.~(\ref{eq:jaa})'' or ``(4.2)''. In multiple-line equations,
the number should be given on the last line.

Displayed equations are to be centered on the page width.  Standard
English letters like x are to appear as $x$ (italicized) in the text
if they are used as mathematical symbols. Punctuation marks are used
at the end of equations as if they appeared directly in the text.

\section{Theorem Environments\label{sfive}}

\begin{theorem}
\label{th1}
Theorems$,$ lemmas$,$ definitions$,$ etc. are set on a separate
paragraph$,$ with extra 1 line space above and below. They are
to be numbered consecutively within each section.
\begin{equation}
\mu(n, t) = \frac{\sum^\infty_{i=1} 1(d_i < t, N(d_i)
= n)}{\int^t_{\sigma=0} 1(N(\sigma) = n)d\sigma}\,.
\label{five1}
\end{equation}

Numbering by sections refers only to the large main sections$,$
i.e. to those with Level 1 headings.  Thus the numbering style in
Sec.~{\upshape\ref{sfive}}
would be labeled {\upshape (\ref{five1}),} {\upshape (5.2),} etc.
\end{theorem}

The theorem environments text can be cross-linked to the body text,
e.g.,\break
Theorem \ref{th1} and Lemma \ref{lem1}.

\begin{lemma}[Sample]
\label{lem1}
Theorems$,$ lemmas$,$ definitions$,$ etc. are set on a separate paragraph$,$
with extra 1 line space above and below. They are to be numbered
consecutively within each section.

Numbering by sections refers only to the large main sections, i.e. to
those with Level 1 headings.  Thus the numbering style in
Sec.~{\upshape\ref{sfive}} would be labeled {\upshape (5.1),}
{\upshape (5.2),} etc.
\end{lemma}

\begin{proof}
Proofs should end with a box.
\end{proof}

\begin{figure}[pb]
\centerline{\includegraphics[width=2.0in]{ijtaff1}}
\vspace*{8pt}
\caption{A schematic illustration of dissociative recombination. The
direct mechanism, 4m$^2_\pi$ is initiated when the
molecular ion S$_{\rm L}$ captures an electron with kinetic energy.}
\end{figure}

\section{Illustrations and Photographs}
Figures are to be inserted in the text nearest their first reference.
Figure placements can be either top or bottom.  Original india ink
drawings of glossy prints are preferred. Please send one set of
originals with copies. If the author requires the publisher to reduce
the figures, ensure that the figures (including letterings and
numbers) are large enough to be clearly seen after reduction. If
photographs are to be used, only black and white ones are acceptable.

Figures are to be sequentially numbered in Arabic numerals. The
caption must be placed below the figure. Typeset in 8 pt roman with
baselineskip of 10~pt. Long captions are to be justified by the
``page-width''.  Use double spacing between a caption and the text
that follows immediately.

Previously published material must be accompanied by written
permission from the author and publisher.

\section{Tables}

Tables should be inserted in the text as close to the point of
reference as possible. Some space should be left above and below
the table.

Tables should be numbered sequentially in the text in Arabic
numerals. Captions are to be centralized above the tables.  Typeset
tables and captions in 8 pt roman with baselineskip of 10 pt. Long
captions are to be justified by the ``table-width''.

\begin{table}[h]
\tbl{Comparison of acoustic for frequencies for piston-cylinder problem.}
{\begin{tabular}{@{}cccc@{}} \toprule
Piston mass & Analytical frequency & TRIA6-$S_1$ model &
\% Error \\
& (Rad/s) & (Rad/s) \\ \colrule
1.0\hphantom{00} & \hphantom{0}281.0 & \hphantom{0}280.81 & 0.07 \\
0.1\hphantom{00} & \hphantom{0}876.0 & \hphantom{0}875.74 & 0.03 \\
0.01\hphantom{0} & 2441.0 & 2441.0\hphantom{0} & 0.0\hphantom{0} \\
0.001 & 4130.0 & 4129.3\hphantom{0} & 0.16\\ \botrule
\end{tabular}}
\begin{tabnote}
Table notes
\end{tabnote}
\begin{tabfootnote}
\tabmark{a} Table footnote A\\
\tabmark{b} Table footnote B
\end{tabfootnote}
\end{table}

If tables need to extend over to a second page, the continuation of
the table should be preceded by a caption, e.g.~``{\it Table~2.}
$(${\it Continued}$)$''. Notes to tables are placed below the final
row of the table and should be flushleft.  Footnotes in tables
shouldbe indicated by superscript lowercase letters and placed beneath
the table.

\section{Running Heads}

Please provide a shortened runninghead (not more than eight words) for
the title of your paper. This will appear on the top right-hand side
of your paper.

\section{Footnotes}

Footnotes should be numbered sequentially in superscript
lowercase roman letters.\footnote{Footnotes should be
typeset in 8 pt roman at the bottom of the page.}

\appendix

\section{Appendices}

Appendices should be used only when absolutely necessary. They
should come\break
before the Acknowledgments. If there is more than one
appendix, number them alphabetically. Number displayed equations
occurring in the Appendix in this way, e.g.~(\ref{appeqn}), (A.2), etc.
\begin{equation}
\mu(n, t) = \frac{\sum^\infty_{i=1} 1(d_i < t, N(d_i)
= n)}{\int^t_{\sigma=0} 1(N(\sigma) = n)d\sigma}\,.
\label{appeqn}
\end{equation}

\section*{Acknowledgments}

This section should come before the References. Funding
information may also be included here.

\section*{References}
References are to be listed in alphabetical order using the last name of the first-named author.
They are to be cited in the text in an author-date format (i.e. Black 1973).
They can be typed after punctuation marks, e.g. ``... in the statement \citep{duf01} ...''
or used directly, e.g. ``see \cite{duf01} for examples.''
For multi-authored papers, please use the `\&' symbol, i.e. \cite{del94}.

Please list the references using the style shown in the following examples:

\def\newsection#1{\section*{\underline{#1}}}

\newsection{Journal articles:}

\noindent R. C. Merton (1974) On the pricing of corporate debt: the risk structure of interest rates,
{\it Journal of Finance} {\bf 29} (2), 449--470.\\[3pt]

\noindent F. Delbaen \& W. Schachermayer (1994) A general version of the fundamental theorem
of asset pricing, {\it Mathematische Annalen} {\bf 300} (1), 463--520.\\[3pt]

\noindent D. C. Brody, L. P. Hughston \& A. Macrina (2008a) Information-based asset pricing,
{\it International Journal of Theoretical and Applied Finance} {\bf 11} (1), 107--142.\\[3pt]

\noindent D. C. Brody, L. P. Hughston \& A. Macrina (2008b) Dam rain and cumulative gain,
{\it Proceedings of the Royal Society A} {\bf 464}, 1801--1822.

\section*{Notes:}
\begin{enumerate}[4.]
\item[1.] List articles alphabetically by the surname of the first-named author.

\item[2.] Author initials should precede the surname.

\item[3.] In the case of multiple listing by the same authors in the same year, the years should be labeled ``a'', ``b'' etc. Thus: (1999a), (1999b), etc.

\item[4.] In the titles of journal articles, the font should be non-italic, and words should be capitalized only on an ``as necessary'' basis.
\end{enumerate}

\newsection{Books:}
\noindent D. Duffie (2001) {\it Dynamic Asset Pricing Theory}, third edition.
Princeton, New Jersey: Princeton University Press.\\[3pt]

\noindent J. Gatheral (2006) {\it The Implied Volatility Surface: A Practitioner's Guide}. New York: Wiley.\\[3pt]

\noindent R. A. Jarrow (2008) {\it Financial Derivatives Pricing: Selected Works of Robert Jarrow}. New
Jersey: World Scientific Publishing Company.\\[3pt]

\noindent I. Karatzas \& S. E. Shreve (1998) {\it Methods of Mathematical Finance}. New York: Springer-Verlag.

\section*{Notes:}
\begin{enumerate}[4.]
\item[1.] In the titles of books, the font should be italic, and all of the ``main'' words should be capitalized.

\item[2.] Generally, avoid abbreviations, except in particular situations where the typesetting can be improved by use of an abbreviation (e.g., to avoid running over a line or a page).

\item[3.] Thus, write ``third edition'' rather than ``3rd ed.''. However, ``ed.'' and ``eds.'' can be used for
``editor'' and ``editors''.
\end{enumerate}

\newsection{Articles appearing in edited books:}

\noindent J. Gatheral, A. Schied \& A. Slynko (2011) Exponential resilience and decay of market impact. In: {\it Econophysics of Order-Driven Markets} (F. Abergel, B. K Chakrabarti, A. Chakraborti \& M. Mitra, eds.), 225--236. Milan: Springer.\\[3pt]

\noindent D. C. Brody, L. P. Hughston \& A. Macrina (2007) Beyond hazard rates: a new framework for credit-risk modeling. In: {\it Advances in Mathematical Finance} (M. C. Fu, R. Jarrow, J. Y. Yen \& R. Elliot, eds.), 231--257. Boston: Birkhauser.

\newsection{PhD theses:}

\noindent A. Macrina (2006) {\it An Information-Based Framework for Asset Pricing: X-Factor Theory and its Applications}. PhD Thesis, Department of Mathematics, King's College London.

\section*{\underline{Articles appearing on ArXiv, but unpublished, or not yet accepted for}\\ \underline{publication:}}

R. A. Jarrow \& M. Larsson (2014) Informational efficiency under short sale constraints. arXiv:1401.1851.

\newsection{Example of reference list at end of article:}

\begin{thebibliography}{9}
\bibitem[Brody {\it et al.}(2007)]{bro07} D. C. Brody, L. P. Hughston \& A. Macrina (2007) Beyond hazard rates: a new framework for credit-risk modeling. In: {\it Advances in Mathematical Finance} (M. C. Fu, R. Jarrow, J. Y. Yen
\& R. Elliot, eds.), 231--257. Boston: Birkhauser.

\bibitem[Brody {\it et al.}(2008a)]{bro08a}D. C. Brody, L. P. Hughston \& A. Macrina (2008a) Information-based asset pricing,
{\it International Journal of Theoretical and Applied Finance} {\bf 11} (1), 107--142.

\bibitem[Brody {\it et al.}(2008b)]{bro08b}D. C. Brody, L. P. Hughston \& A. Macrina (2008b) Dam rain and cumulative gain,
{\it Proceedings of the Royal Society A} {\bf 464}, 1801--1822.

\bibitem[Delbaen \& Schachermayer(1994)]{del94} F. Delbaen \& W. Schachermayer (1994) A general version of the fundamental theorem of asset pricing, {\it Mathematische Annalen} {\bf 300} (1), 463--520.

\bibitem[Duffie(2001)]{duf01} D. Duffie (2001) {\it Dynamic Asset Pricing Theory}, third edition. Princeton, New Jersey: Princeton University Press.

\bibitem[Gatheral(2006)]{gat06} J. Gatheral (2006) {\it The Implied Volatility Surface: A Practitioner's Guide}. New York: Wiley.

\bibitem[Gatheral {\it et al.}(2011)]{gat11} J. Gatheral, A. Schied \& A. Slynko (2011) Exponential resilience and decay of market impact. In: {\it Econophysics of Order-Driven Markets} (F. Abergel, B. K Chakrabarti, A. Chakraborti \& M. Mitra, eds.), 225--236. Milan: Springer.

\bibitem[Jarrow(2008)]{jar08} R. A. Jarrow (2008) {\it Financial Derivatives Pricing: Selected Works of Robert Jarrow}. New
Jersey: World Scientific Publishing Company.

\bibitem[Jarrow & Larsson(2014)]{jar14} R. A. Jarrow \& M. Larsson (2014) Informational efficiency under short sale constraints. arXiv:1401.1851.

\bibitem[Karatzas \& Shreve(1998)]{kar98} I. Karatzas \& S. E. Shreve (1998) {\it Methods of Mathematical Finance}. New York: Springer- Verlag.

\bibitem[Macrina(2006)]{mac06} A. Macrina (2006) {\it An Information-Based Framework for Asset Pricing: X-Factor Theory and its Applications}. PhD Thesis, Department of Mathematics, King's College London.

\bibitem[Merton(1974)]{mer74} R. C. Merton (1974) On the pricing of corporate debt: the risk structure of interest rates,
{\it Journal of Finance} {\bf 29} (2), 449--470.

\bibitem[Madan & Schoutens(2011)]{mad11} D. Madan \& W. Schoutens (2011) Conic finance and the corporate balance sheet, {\it International Journal of Theoretical and Applied Finance} {\bf 14} (5) 587--610. Reprinted (2013) in: {\it Finance at Fields} (M. R. Grasselli \& L. P. Hughston, eds.), 451--474.
\end{thebibliography}

\section*{Note:}

References are to be listed in alphabetical order using the last name of the first-named author.

\newsection{Examples of citations:}

\begin{enumerate}[(f)]
\item[(a)] This result was first established by \cite{mer74}.
\item[(b)] This result is discussed in \cite{mer74}.
\item[(c)] The arbitrage-free approach to interest rate theory has a long history \citep{mer74}.
\item[(d)] The theory of arbitrage free pricing has been greatly advanced in recent years \citep{del94,kar98,duf01}.
\item[(e)] See \cite{del94,kar98,duf01}, for further information.
\item[(f)] The role of information in finance is generally appreciated \citep{bro07,bro08a,bro08b}.
\item[(g)] For more details about the role of information, see \cite{bro07,bro08a,bro08b}.
\end{enumerate}
\end{document} 